\subsection{Hyperparameter learning and the PES acquisition function}
\label{sec:hyper-learning}

We now show how the previous approximations are integrated to compute the acquisition function used by 
predictive entropy search (PES).
This acquisition function performs a formal treatment of the hyperparameters. Let $\vpsi$ denote
a vector of hyperparameters which includes any kernel parameters as well as the
noise variance $\sigma^2$. Let $p(\vpsi|\D_n)\propto p(\vpsi)
\,p(\D_n|\vpsi)$ denote the posterior distribution over these
parameters where $p(\vpsi)$ is a hyperprior  and $p(\D_n|\vpsi)$ is the GP marginal likelihood.
For a fully Bayesian treatment of $\vpsi$ we must marginalize the acquisition
function (\ref{eq:newObjective}) with respect to this posterior. 
The corresponding integral has no analytic expression and must be approximated using Monte
Carlo. This approach is also taken in \cite{Snoek:2012}.

We draw $M$ samples $\{\vpsi\i\}$ from $p(\vpsi|\D_n)$ using
slice sampling~\cite{Vanhatalo:2012}. Let $\xopt\i$ denote a sampled global maximizer
drawn from $p(\xopt|\D_n,\vpsi\i)$ as described in Section \ref{sec:sampling}.  Furthermore, let $v_n\i(\x)$ and $v_n\i(\x|\xopt\i)$ denote
the predictive variances computed as described in
Section~\ref{sec:predictiveEntropy} when the model hyperparameters are fixed to $\vpsi\i$. 
We then write the marginalized acquisition function as
\begin{align}
    \alpha_n(\x) 
    &=
    \textstyle
    \frac{1}{M}\sum_{i=1}^M 
    \left\{ 0.5\log[v_n\i(\x)+\sigma^2] - 0.5\log[v_n\i(\x|\xopt\i)+\sigma^2]\right\}
    \,.\label{eq:marginalizedObjective}
\end{align}
Note that PES is effectively marginalizing the original acquisition function
(\ref{eq:originalObjective}) over $p(\vpsi|\D_n)$. This is a significant
advantage with respect to other methods that optimize the same
information-theoretic acquisition function but do not marginalize over
the hyper-parameters. For example, the approach of
\cite{HennigSchuler:2012} approximates
(\ref{eq:originalObjective}) only for fixed $\vpsi$. The resulting approximation is
computationally very expensive and recomputing it to average over multiple
samples from $p(\vpsi|\D_n)$ is infeasible in practice.  

Algorithm~\ref{alg:pes} shows pseudo-code for computing the PES acquisition
function. Note that most of the computations necessary for evaluating (\ref{eq:marginalizedObjective}) can be done
independently of the input $\x$, as noted in the pseudo-code. This initial cost
is dominated by a matrix inversion necessary to pre-compute $\vV$ for each
hyperparameter sample. The resulting complexity is $\calO[M(n+d + d(d-1)/2)^3]$.
This cost can be reduced to $\calO[M(n+d)^3]$ by ignoring the 
derivative observations imposed on $\upper[\nabla^2 f(\xopt)]$ by constraint C1.1.
Nevertheless, in the problems that we consider $d$ is very small (less than 20).
After these precomputations are done, the evaluation of (\ref{eq:marginalizedObjective})
is $\calO[M(n+d + d(d-1)/2)]$.

%Note also that the
%computational complexity of PES is dominated by the matrix inversion in
%(\ref{eq:mpredvpred}).  This operation has to be performed for each $\xopt\i$
%and as a result the cost of each iteration of PES is $\calO[M(n + d +
%d(d-1)/2)^3]$.

